\documentclass[UTF8,11pt]{ctexart}

\usepackage[a4paper,margin=1in]{geometry}
\usepackage{amsmath,amssymb,bm}
\usepackage{booktabs}
\usepackage[hidelinks]{hyperref}

\title{固定层级随机插值速度场的一步生成推导}
\author{}
\date{\today}

\newcommand{\E}{\mathbb{E}}
\newcommand{\R}{\mathbb{R}}
\newcommand{\calN}{\mathcal{N}}
\newcommand{\dd}{\,\mathrm{d}}

\begin{document}

\maketitle

\section{一般形式}

随机插值的一般形式可以写为
\begin{equation}
I_\tau
=
a_\tau X_0
+
b_\tau X_1
+
\gamma_\tau \varepsilon,
\qquad
\varepsilon \sim \calN(0,I_D),
\qquad
\tau\in(0,1).
\end{equation}

其中 \(X_0\sim p_0\) 是 base endpoint，\(X_1\sim p_{\mathrm{data}}\) 是数据 endpoint。
速度场学习的目标不是 score，而是条件平均速度：
\begin{equation}
v_\tau(y)
=
\E
\left[
\dot a_\tau X_0
+
\dot b_\tau X_1
+
\dot \gamma_\tau \varepsilon
\mid
I_\tau=y
\right].
\end{equation}

如果希望最终做一步生成，引入非可逆低维传输映射
\begin{equation}
X_1^\theta=h_\theta(U),
\qquad
U\sim \pi,
\qquad
h_\theta:\R^d\to\R^D,
\end{equation}
采样时直接输出
\begin{equation}
x_{\mathrm{gen}}=h_\theta(u),
\qquad
u\sim\pi.
\end{equation}

\section{情况一：无噪声线性插值}

取
\begin{equation}
a_\tau = 1-\tau,
\qquad
b_\tau = \tau,
\qquad
\gamma_\tau = 0.
\end{equation}
于是
\begin{equation}
X_\tau
=
(1-\tau)X_0+\tau X_1.
\end{equation}

因此 真实速度场是
\begin{equation}
\dot X_\tau
=
X_1-X_0.
\end{equation}


一步生成器诱导的插值为
\begin{equation}
X_{\theta,\tau}
=
(1-\tau)X_0+\tau h_\theta(U).
\end{equation}
模型速度场为
\begin{equation}
   \begin{aligned}
    v_{\mathrm{data},\tau}(x_\tau)&=\E\left[h_\theta(U)-X_0\mid X_\tau=x_\tau\right]\\
    &=\E\left[h_\theta(U)-\frac{X_\tau-\tau h_\theta(U)}{1-\tau}\mid X_\tau=x_\tau\right]\\
    &=\frac{1}{1-\tau}\E\left[h_\theta(U)\mid X_\tau=x_\tau\right]-\frac{1}{1-\tau}x_\tau.
   \end{aligned}
\end{equation}

对应的 velocity matching 目标可写为
\begin{equation}
\mathcal{L}_{\mathrm{lin}}(\theta)
=
\E
\left[
\left\|
v_{\theta,\tau}(I_\tau)
-
(X_1-X_0)
\right\|_2^2
\right].
\end{equation}

\subsection{Anchor 近似}

由于 \(\gamma_\tau=0\)，条件分布是退化的，并不会自然产生 Gaussian likelihood 权重。
工程上可以加入一个很小的 kernel bandwidth \(\eta>0\)，用核平滑近似条件分布。
给定 anchors \((u_k)_{k=1}^K\)，定义
\begin{equation}
z_k
=
-
\frac{
\left\|
y-(1-\tau)x_{0,k}-\tau h_\theta(u_k)
\right\|_2^2
}{
2\eta^2
},
\end{equation}
\begin{equation}
w_k
=
\frac{\exp(z_k)}
{\sum_{j=1}^K\exp(z_j)}.
\end{equation}
于是
\begin{equation}
v_{\theta,K,\tau}(y)
=
\sum_{k=1}^{K}
w_k
\left[
h_\theta(u_k)-x_{0,k}
\right].
\end{equation}

这个无噪声版本非常接近 rectified flow 或 flow matching，但若要使用 MAGT-style anchor 推导，通常需要额外的核平滑或显式速度网络。

\section{情况二：线性插值加中间随机扰动}

取
\begin{equation}
a_\tau=1-\tau,
\qquad
b_\tau=\tau,
\qquad
\gamma_\tau=\sqrt{2\tau(1-\tau)}.
\end{equation}
于是
\begin{equation}
I_\tau
=
(1-\tau)X_0
+
\tau X_1
+
\sqrt{2\tau(1-\tau)}\,\varepsilon.
\end{equation}

记
\begin{equation}
\gamma_\tau=\sqrt{2\tau(1-\tau)}.
\end{equation}
则
\begin{equation}
\dot a_\tau=-1,
\qquad
\dot b_\tau=1,
\qquad
\dot\gamma_\tau
=
\frac{1-2\tau}
{\sqrt{2\tau(1-\tau)}}.
\end{equation}

因此 sample-wise velocity 为
\begin{equation}
\dot I_\tau
=
X_1-X_0
+
\frac{1-2\tau}
{\sqrt{2\tau(1-\tau)}}\varepsilon.
\end{equation}

真实速度场为
\begin{equation}
v_{\mathrm{data},\tau}(y)
=
\E
\left[
X_1-X_0
+
\frac{1-2\tau}
{\sqrt{2\tau(1-\tau)}}\varepsilon
\mid
I_\tau=y
\right].
\end{equation}

由
\begin{equation}
\varepsilon
=
\frac{
y-(1-\tau)X_0-\tau X_1
}{
\sqrt{2\tau(1-\tau)}
},
\end{equation}
代入可得
\begin{equation}
\dot I_\tau
=
X_1-X_0
+
\frac{1-2\tau}
{2\tau(1-\tau)}
\left[
y-(1-\tau)X_0-\tau X_1
\right].
\end{equation}

整理后得到
\begin{equation}
\dot I_\tau
=
\frac{1-2\tau}
{2\tau(1-\tau)}y
-
\frac{1}{2\tau}X_0
+
\frac{1}{2(1-\tau)}X_1.
\end{equation}

因此真实速度场也可写为
\begin{equation}
v_{\mathrm{data},\tau}(y)
=
\frac{1-2\tau}
{2\tau(1-\tau)}y
-
\frac{1}{2\tau}
\E[X_0\mid I_\tau=y]
+
\frac{1}{2(1-\tau)}
\E[X_1\mid I_\tau=y].
\end{equation}

\section{一步生成器诱导的速度场}

令
\begin{equation}
X_1^\theta=h_\theta(U),
\qquad
U\sim \pi.
\end{equation}
模型插值为
\begin{equation}
I_{\theta,\tau}
=
(1-\tau)X_0
+
\tau h_\theta(U)
+
\sqrt{2\tau(1-\tau)}\,\varepsilon.
\end{equation}

模型速度场为
\begin{equation}
v_{\theta,\tau}(y)
=
\E
\left[
h_\theta(U)-X_0
+
\dot\gamma_\tau\varepsilon
\mid
I_{\theta,\tau}=y
\right].
\end{equation}

\subsection{Anchor 权重}

用 anchors \((x_{0,k},u_k)_{k=1}^K\) 近似条件期望，记
\begin{equation}
h_k=h_\theta(u_k).
\end{equation}
因为噪声方差为 \(\gamma_\tau^2=2\tau(1-\tau)\)，Gaussian likelihood 对应的 log-weight 为
\begin{equation}
z_k
=
-
\frac{
\left\|
y-(1-\tau)x_{0,k}-\tau h_k
\right\|_2^2
}{
4\tau(1-\tau)
}
+
\log p_0(x_{0,k})
+
\log \pi(u_k)
-
\log q(x_{0,k},u_k\mid y).
\end{equation}
归一化权重为
\begin{equation}
w_k
=
\frac{\exp(z_k)}
{\sum_{j=1}^{K}\exp(z_j)}.
\end{equation}

对每个 anchor，有
\begin{equation}
\varepsilon_k(y)
=
\frac{
y-(1-\tau)x_{0,k}-\tau h_k
}{
\sqrt{2\tau(1-\tau)}
}.
\end{equation}

于是模型速度场近似为
\begin{equation}
v_{\theta,K,\tau}(y)
=
\sum_{k=1}^{K}
w_k
\left[
h_k-x_{0,k}
+
\frac{1-2\tau}
{\sqrt{2\tau(1-\tau)}}
\varepsilon_k(y)
\right].
\end{equation}

也可以整理成更紧凑的形式：
\begin{equation}
v_{\theta,K,\tau}(y)
=
\frac{1-2\tau}
{2\tau(1-\tau)}y
+
\sum_{k=1}^{K}
w_k
\left[
-
\frac{1}{2\tau}x_{0,k}
+
\frac{1}{2(1-\tau)}h_k
\right].
\end{equation}

\section{训练目标}

带中间随机扰动时，真实 sample-wise target 为
\begin{equation}
v_{\mathrm{target}}
=
X_1-X_0
+
\frac{1-2\tau}
{\sqrt{2\tau(1-\tau)}}\varepsilon.
\end{equation}

等价地，也可以写为
\begin{equation}
v_{\mathrm{target}}
=
\frac{1-2\tau}
{2\tau(1-\tau)}I_\tau
-
\frac{1}{2\tau}X_0
+
\frac{1}{2(1-\tau)}X_1.
\end{equation}

因此训练目标为
\begin{equation}
\mathcal{L}_{\mathrm{SI}}(\theta)
=
\E
\left[
\left\|
v_{\theta,K,\tau}(I_\tau)
-
v_{\mathrm{target}}
\right\|_2^2
\right].
\end{equation}

其中
\begin{equation}
I_\tau
=
(1-\tau)X_0
+
\tau X_1
+
\sqrt{2\tau(1-\tau)}\,\varepsilon.
\end{equation}

\section{实现备注}

无噪声线性插值
\[
I_\tau=(1-\tau)X_0+\tau X_1
\]
形式简单，但条件分布退化；若想使用 anchor posterior 的权重形式，需要人为加入 kernel smoothing。

带随机扰动的插值
\[
I_\tau=(1-\tau)X_0+\tau X_1+\sqrt{2\tau(1-\tau)}\,\varepsilon
\]
更适合 MAGT-style 推导，因为它自然给出 Gaussian likelihood 权重，并可通过 anchors 近似条件速度场。

最终的一步生成仍然不需要积分速度场，而是直接使用
\begin{equation}
x_{\mathrm{gen}}=h_\theta(u),
\qquad
u\sim\pi.
\end{equation}

\end{document}
